% !TEX root = goalmisgen-slides.tex
% goalmisgen-slides.tex - Specification Gaming & Goal Misgeneralisation (ILIAD Intensive).
% Written from the ARENA [2.6] master (master_2_6.py). Madrid/Beamer via arena_beamer.sty.
% Build:  present -> pdflatex goalmisgen-slides.tex
%         handout -> pdflatex "\def\HANDOUT{}% !TEX root = goalmisgen-slides.tex
% goalmisgen-slides.tex - Specification Gaming & Goal Misgeneralisation (ILIAD Intensive).
% Written from the ARENA [2.6] master (master_2_6.py). Madrid/Beamer via arena_beamer.sty.
% Build:  present -> pdflatex goalmisgen-slides.tex
%         handout -> pdflatex "\def\HANDOUT{}% !TEX root = goalmisgen-slides.tex
% goalmisgen-slides.tex - Specification Gaming & Goal Misgeneralisation (ILIAD Intensive).
% Written from the ARENA [2.6] master (master_2_6.py). Madrid/Beamer via arena_beamer.sty.
% Build:  present -> pdflatex goalmisgen-slides.tex
%         handout -> pdflatex "\def\HANDOUT{}% !TEX root = goalmisgen-slides.tex
% goalmisgen-slides.tex - Specification Gaming & Goal Misgeneralisation (ILIAD Intensive).
% Written from the ARENA [2.6] master (master_2_6.py). Madrid/Beamer via arena_beamer.sty.
% Build:  present -> pdflatex goalmisgen-slides.tex
%         handout -> pdflatex "\def\HANDOUT{}\input{goalmisgen-slides}"
\input{arena_beamer.sty}

% deck-local macros
\newcommand{\Ra}{\mathcal{R}}
\newcommand{\St}{\mathcal{S}}
\newcommand{\Ac}{\mathcal{A}}

\begin{document}

\title[{[2.6] Goal Misgeneralisation}]{{[2.6] Specification Gaming \& Goal Misgeneralisation}}
\author[David Quarel]{David Quarel}
\institute[ILIAD]{ILIAD Intensive}
\date[\today]{\today}
\frame{\titlepage}

% -------------------------------------------------------------------
\begin{frame}{Two ways reward goes wrong}
We tell an RL agent what we want through a \textbf{reward function}. Two distinct failure modes:
\begin{enumerate}
  \item[] \pause
  \item \textbf{Goal misspecification} (reward hacking): the designer specified
  a reward function that incentivises a behaviour that does not match what the designer intended.
  The reward was optimised, but the desired goal was not internalised by the agent.
  \pause
  \item \textbf{Goal misgeneralisation}: the reward function matches the intended goal,
  but the agent learns a goal that is perfectly correlated with the intended goal 
  \emph{on the training distribution}, but which diverges on out-of-distribution environments.
  The agent remains capable, but capably pursues the wrong thing.
\end{enumerate}
\pause
\vfill
Even with a \emph{perfectly specified} reward, it is possible for the agent 
to still internalise the wrong goal.
\end{frame}

% -------------------------------------------------------------------
\begin{frame}{The pottery shop}
\begin{columns}
\begin{column}{0.5\textwidth}
A grid world: a robot shares the floor with fragile \textbf{urns} and piles of \textbf{shards}, 
plus a \textbf{bin}. We want it to tidy shards into the bin without smashing urns.
\begin{itemize}
  \item Actions: move around, \texttt{PICKUP}, \texttt{PUTDOWN}.
  \item Walking into an urn breaks it into shards.
  \item Shards are disposed of by picking them up, and dropping them in the bin.
  \item Play this environment in your browser: \href{https://iliad-team.github.io/iliad-intensive-D.2/play.html}{https://iliad-team.github.io/iliad-intensive-D.2/play.html}
\end{itemize}
\end{column}
\begin{column}{0.5\textwidth}
\centering
\imgorplaceholder{0.75\textheight}{pottery_shop}{pottery shop env}\\[1ex]
\end{column}
\end{columns}
\end{frame}

% -------------------------------------------------------------------
\begin{frame}{Reward functions as the interface}

\begin{block}{The reward hypothesis}
  ``That all of what we mean by goals and purposes can be well thought of as maximization of the expected value of the cumulative sum of a received scalar signal (reward).'' -- Richard Sutton
\end{block}
Assuming this hypothesis, it should be possible (at least in principle) to encourage 
any sort of behaviour we want by finding just the right reward function.
\pause
\vfill
We want to train the agent to pick shards up and carry them to the bin.
A first attempt, \texttt{reward1}: 
\begin{enumerate}
\item $+1$ for picking up shards, 
\item $+1$ for dropping them in the bin.
\end{enumerate}
What behaviour does this reinforce in practice?
\end{frame}

% -------------------------------------------------------------------
\begin{frame}{Specification gaming: the agent farms reward}
\begin{columns}
\begin{column}{0.43\textwidth}
Trained on \texttt{reward1}, the agent discovers a loophole:
\begin{itemize}
  \item each \texttt{PICKUP} is $+1$,
  \item dropping on the floor costs nothing,
  \item so it picks up and drops the same pile over and over.
\end{itemize}
\pause
We probe behaviour with extra reward functions: \texttt{reward\_drop} ($+1$ per floor-drop). Its return tracks \texttt{reward1} almost exactly.
\end{column}
\begin{column}{0.43\textwidth}
\centering
\imgorplaceholder{\linewidth}{hist_specgaming}{spec-gaming histogram}
\end{column}
\end{columns}
\end{frame}

% -------------------------------------------------------------------
\begin{frame}{Fix 1: potential shaping (the un-hackable hint)}
We still want to \emph{hint} that holding shards is good, without creating a loophole. \\
Removing the reward from picking up shards makes the task very hard to learn.
\vfill
\textbf{Potential shaping} (Ng, Harada \& Russell, 1999)~\citev{Ng:99shaping}: pick a bounded $\Phi: \St \to \mathbb{R}$ and define a new reward function
\[
  R^\Phi(s,a,s') = R(s,a,s') + \gamma\,\Phi(s') - \Phi(s).
\]
\pause
Over a trajectory the $\Phi$ terms \textbf{telescope}: total return changes only by a constant $-\Phi(s_0)$, so the ordering over policies (and the optimum) is unchanged. It is the \emph{only} shaping with this guarantee.
\pause
\vfill
Here $\Phi(s) = 0.5$ while holding shards. A pickup advances $+\gamma\Phi$; a floor-drop gives it back. A \textbf{pickup-then-drop cycle nets exactly zero}: the loophole is closed, while also encouraging the agent
to try picking up shards in the first place.
\end{frame}

% -------------------------------------------------------------------
\begin{frame}{Why shaping preserves the optimum}
Write the return from time $t$ with $r_k := R(s_{k-1},a_{k-1},s_k)$:
\[
  G_t = r_{t+1} + \gamma\, r_{t+2} + \gamma^2 r_{t+3} + \cdots
      = \sum_{k=t+1}^{\infty} \gamma^{\,k-(t+1)}\, r_k.
\]
\pause
The shaped reward just adds a potential difference, $r_k^\Phi = r_k + \gamma\,\Phi(s_k) - \Phi(s_{k-1})$, so
\[
  G_t^\Phi = \sum_{k=t+1}^{\infty} \gamma^{\,k-(t+1)}\, r_k^\Phi
           = G_t + \sum_{k=t+1}^{\infty} \gamma^{\,k-(t+1)}\big(\gamma\,\Phi(s_k) - \Phi(s_{k-1})\big).
\]
\pause
Expanding the shaping sum, consecutive terms \textbf{telescope}:
\[
  \big(\gamma\Phi(s_{t+1}) - \Phi(s_t)\big)
  + \gamma\big(\gamma\Phi(s_{t+2}) - \Phi(s_{t+1})\big)
  + \cdots = -\Phi(s_t).
\]
\pause
Hence $G_t^\Phi = G_t - \Phi(s_t)$. At $t=0$, $G_0^\Phi = G_0 - \Phi(s_0)$ differs from $G_0$ by a constant depending \emph{only on the start state}, not on the policy: so $\arg\max_\pi G_0^\Phi = \arg\max_\pi G_0$ and the optimum is unchanged.
\end{frame}

% % -------------------------------------------------------------------
% \begin{frame}{Fix 2: penalise smashing urns}

% % \pause
% % \[
% %   \texttt{reward2} = \underbrace{\texttt{reward\_shaped}}_{\Phi=0.5,\ \text{bin }+1} + \underbrace{\texttt{reward\_no\_break}}_{-2 \text{ per smash}}
% % \]
% \end{frame}

% -------------------------------------------------------------------
\begin{frame}[t]{Specification gaming fixed?}
\begin{columns}[T]
\begin{column}{0.5\textwidth}
Shaping stops farming, but nothing yet discourages breaking urns to manufacture shards.
\begin{itemize}
  \item Add a $-2$ penalty for each urn smashed.
  \item Binning the resulting shard recovers $\approx +1$ (discounted), against a $-2$ penalty: 
  so breaking-then-cleaning is never worth it.
\end{itemize}
Retrained on \texttt{reward2}, the probes confirm the fix:
\begin{itemize}
  \item \texttt{reward\_drop} (farming) $\approx 0$,
  \item \texttt{reward\_break} (smashing) $\approx 0$,
  \item \texttt{reward\_bin} $> 0$: it actually cleans.
\end{itemize}


\end{column}
\begin{column}{0.5\textwidth}
\centering
\imgorplaceholder{0.55\linewidth}{hist_fixed}{fixed histogram}\\[1ex]
\raggedright\footnotesize

\end{column}
\end{columns}
\end{frame}


% -------------------------------------------------------------------
\begin{frame}{Motivating the potential function}
\begin{itemize}
  \item Why bother with the potential function at all? +1 reward for putting shards in
  the bin, and -2 for smashing urns, also is aligned with the intended goal.
  \pause
  \item But the agent still has to learn from experience: the policy does not
  teleport to whatever maximises the sum of rewards.
  \pause
  \item Leaving the potential function out robs the agent of useful feedback (that picking up
  and walking around with shards is instrumentally useful for putting them in the bin).
\end{itemize}
\end{frame}

\begin{frame}{}
  \begin{center}
    \imgorplaceholder{0.8\textwidth}{potential_train}{}
    \end{center}
\end{frame}


% -------------------------------------------------------------------
\begin{frame}{Even perfectly specified reward isn't enough}
A policy trained on a \emph{single} fixed layout can just memorise a trajectory: it never has to look at where the shards are.
\pause
\vfill
So we train on \textbf{procedurally generated} layouts: random robot and item positions every episode. \\
As a toy example of insufficent training environment diversity, suppose that the bin is always in the top-left corner.
%\vspace{-0.6cm}
\begin{center}
\imgorplaceholder{0.75\linewidth}{generate_grid}{}
\end{center}
\end{frame}

% -------------------------------------------------------------------
\begin{frame}{Goal misgeneralisation: defining it}
\textbf{Langosco et al. (2022)}~\citev{Langosco:22goalmisgen}: an agent trained to maximise $R$ is deployed 
under distribution shift. \emph{Goal misgeneralisation} occurs if it now gets \textbf{low reward} 
(according to $R$)because it keeps acting as if it was trained to optimise a \emph{different} $R' \neq R$.
\begin{itemize}
  \item $R$ = \textbf{intended objective} (carry shards to the bin),
  \item $R'$ = \textbf{behavioural objective} (what it actually learned).
\end{itemize}
\pause
During training, ``go to the bin'' and ``go to the top-left corner'' were \textbf{indistinguishable}. 
The simpler one can win, and the model has internalised the wrong goal.
\end{frame}

% -------------------------------------------------------------------
\begin{frame}{Out of distribution: move the bin}
\begin{columns}
\begin{column}{0.5\textwidth}
At test time, put the bin somewhere else (e.g. the top-right corner).
\pause
The agent still capably picks up shards and carries them \textbf{to the top-left corner}, where it drops them on the floor, ignoring the actual bin.
\pause
Defining a proxy reward function $R'$ that rewards dropping shards at the top left corner,
we can probe the agent's behaviour on the shifted layout.
\end{column}
\begin{column}{0.5\textwidth}
\centering
\imgorplaceholder{\linewidth}{env_shift}{shifted env (bin moved)}
\end{column}
\end{columns}
\end{frame}

% -------------------------------------------------------------------
\begin{frame}{Goal misgeneralisation, measured}
\begin{columns}
\begin{column}{0.5\textwidth}
On shifted layouts, the agent:
\begin{itemize}
  \item scores \emph{much lower} on $R$ (putting shards in the bin) than in-distribution,
  \item but high on $R'$ (dropping shards at the top left corner).
\end{itemize}
\pause
It is still competently optimising a goal, just not what we trained for. 
There is no black-box way to tell apart an agent that learned $R$ from one that learned $R'$
as the behaviour is the same on in-distribution layouts.
\end{column}
\begin{column}{0.5\textwidth}
\centering
\imgorplaceholder{0.8\linewidth}{hist_misgen}{misgen histogram}
\end{column}
\end{columns}
\end{frame}

% -------------------------------------------------------------------
\begin{frame}[t]{Fixing it: broaden the training distribution}
\begin{columns}[T]
\begin{column}{0.5\textwidth}  
\begin{itemize}
\item Randomise the bin position too.
\item Now ``corner'' no longer correlates with ``bin'',
so the agent is forced to learn the intended goal $R$.
\pause
\item Result: The training distribution now matches the test distribution,
so the agent is no longer misgeneralising.
\pause
\item \textbf{Caveat:} In practice you often cannot enumerate the shift in advance,
so this fix is not always available. Mitigating goal misgeneralisation without it is an open problem.
\end{itemize}
\end{column}
\begin{column}{0.5\textwidth}
\centering
\imgorplaceholder{0.8\linewidth}{hist_misgen_fixed}{fixed misgen histogram}
\end{column}
\end{columns}
\end{frame}

% -------------------------------------------------------------------
% \begin{frame}{Aside: everyone has a price}
% \begin{columns}
% \begin{column}{0.52\textwidth}
% With $-2$ for smashing, is breaking an urn \emph{ever} optimal?
% \pause
% \begin{itemize}
%   \item Wall off a pile of shards behind urns, agent has to walk the long way around.
%   \item Beaking an urn at the top opens a shortcut between the shards and the bin.
%   \item Break one urn ($-2$): reach the pile and bin it, return $> 0$.
% \end{itemize}
% \pause
% Optimal policy \textbf{breaks the wall}. But trained PPO often \emph{never discovers it}: it must pay an immediate, salient $-2$ before any reward is reachable. Incentivised $\neq$ learned.
% \end{column}
% \begin{column}{0.48\textwidth}
% \centering
% \imgorplaceholder{\linewidth}{everyone_has_a_price}{impassable wall}
% \end{column}
% \end{columns}
% \end{frame}

% -------------------------------------------------------------------
\begin{frame}{Summary}
\begin{itemize}
  \item \textbf{Specification gaming / outer alignment:} if a behaviour scores higher return than the intended one, the policy will pursue it. Probe with diagnostic rewards; fix with potential shaping (un-hackable) and targeted penalties.
  \pause
  \item \textbf{Goal misgeneralisation / inner alignment:} out of distribution, behaviour is set by the inductive biases, not by what the reward \emph{would} have rewarded. A simpler goal that perfectly correlated during training is what gets pursued when goals diverge.
  \pause
  \item Broadening the training distribution helps, when you can; and ``reward-maximising'' is necessary but not sufficient for ``learned''.
\end{itemize}
\end{frame}

% -------------------------------------------------------------------
\begin{frame}[allowframebreaks]{References}
\footnotesize
\bibliographystyle{alphaurl}
\bibliography{biblo}
\nocite{Langosco:22goalmisgen, Shah:22goalmisgen, Ng:99shaping}
\end{frame}

\end{document}
"
\input{arena_beamer.sty}

% deck-local macros
\newcommand{\Ra}{\mathcal{R}}
\newcommand{\St}{\mathcal{S}}
\newcommand{\Ac}{\mathcal{A}}

\begin{document}

\title[{[2.6] Goal Misgeneralisation}]{{[2.6] Specification Gaming \& Goal Misgeneralisation}}
\author[David Quarel]{David Quarel}
\institute[ILIAD]{ILIAD Intensive}
\date[\today]{\today}
\frame{\titlepage}

% -------------------------------------------------------------------
\begin{frame}{Two ways reward goes wrong}
We tell an RL agent what we want through a \textbf{reward function}. Two distinct failure modes:
\begin{enumerate}
  \item[] \pause
  \item \textbf{Goal misspecification} (reward hacking): the designer specified
  a reward function that incentivises a behaviour that does not match what the designer intended.
  The reward was optimised, but the desired goal was not internalised by the agent.
  \pause
  \item \textbf{Goal misgeneralisation}: the reward function matches the intended goal,
  but the agent learns a goal that is perfectly correlated with the intended goal 
  \emph{on the training distribution}, but which diverges on out-of-distribution environments.
  The agent remains capable, but capably pursues the wrong thing.
\end{enumerate}
\pause
\vfill
Even with a \emph{perfectly specified} reward, it is possible for the agent 
to still internalise the wrong goal.
\end{frame}

% -------------------------------------------------------------------
\begin{frame}{The pottery shop}
\begin{columns}
\begin{column}{0.5\textwidth}
A grid world: a robot shares the floor with fragile \textbf{urns} and piles of \textbf{shards}, 
plus a \textbf{bin}. We want it to tidy shards into the bin without smashing urns.
\begin{itemize}
  \item Actions: move around, \texttt{PICKUP}, \texttt{PUTDOWN}.
  \item Walking into an urn breaks it into shards.
  \item Shards are disposed of by picking them up, and dropping them in the bin.
  \item Play this environment in your browser: \href{https://iliad-team.github.io/iliad-intensive-D.2/play.html}{https://iliad-team.github.io/iliad-intensive-D.2/play.html}
\end{itemize}
\end{column}
\begin{column}{0.5\textwidth}
\centering
\imgorplaceholder{0.75\textheight}{pottery_shop}{pottery shop env}\\[1ex]
\end{column}
\end{columns}
\end{frame}

% -------------------------------------------------------------------
\begin{frame}{Reward functions as the interface}

\begin{block}{The reward hypothesis}
  ``That all of what we mean by goals and purposes can be well thought of as maximization of the expected value of the cumulative sum of a received scalar signal (reward).'' -- Richard Sutton
\end{block}
Assuming this hypothesis, it should be possible (at least in principle) to encourage 
any sort of behaviour we want by finding just the right reward function.
\pause
\vfill
We want to train the agent to pick shards up and carry them to the bin.
A first attempt, \texttt{reward1}: 
\begin{enumerate}
\item $+1$ for picking up shards, 
\item $+1$ for dropping them in the bin.
\end{enumerate}
What behaviour does this reinforce in practice?
\end{frame}

% -------------------------------------------------------------------
\begin{frame}{Specification gaming: the agent farms reward}
\begin{columns}
\begin{column}{0.43\textwidth}
Trained on \texttt{reward1}, the agent discovers a loophole:
\begin{itemize}
  \item each \texttt{PICKUP} is $+1$,
  \item dropping on the floor costs nothing,
  \item so it picks up and drops the same pile over and over.
\end{itemize}
\pause
We probe behaviour with extra reward functions: \texttt{reward\_drop} ($+1$ per floor-drop). Its return tracks \texttt{reward1} almost exactly.
\end{column}
\begin{column}{0.43\textwidth}
\centering
\imgorplaceholder{\linewidth}{hist_specgaming}{spec-gaming histogram}
\end{column}
\end{columns}
\end{frame}

% -------------------------------------------------------------------
\begin{frame}{Fix 1: potential shaping (the un-hackable hint)}
We still want to \emph{hint} that holding shards is good, without creating a loophole. \\
Removing the reward from picking up shards makes the task very hard to learn.
\vfill
\textbf{Potential shaping} (Ng, Harada \& Russell, 1999)~\citev{Ng:99shaping}: pick a bounded $\Phi: \St \to \mathbb{R}$ and define a new reward function
\[
  R^\Phi(s,a,s') = R(s,a,s') + \gamma\,\Phi(s') - \Phi(s).
\]
\pause
Over a trajectory the $\Phi$ terms \textbf{telescope}: total return changes only by a constant $-\Phi(s_0)$, so the ordering over policies (and the optimum) is unchanged. It is the \emph{only} shaping with this guarantee.
\pause
\vfill
Here $\Phi(s) = 0.5$ while holding shards. A pickup advances $+\gamma\Phi$; a floor-drop gives it back. A \textbf{pickup-then-drop cycle nets exactly zero}: the loophole is closed, while also encouraging the agent
to try picking up shards in the first place.
\end{frame}

% -------------------------------------------------------------------
\begin{frame}{Why shaping preserves the optimum}
Write the return from time $t$ with $r_k := R(s_{k-1},a_{k-1},s_k)$:
\[
  G_t = r_{t+1} + \gamma\, r_{t+2} + \gamma^2 r_{t+3} + \cdots
      = \sum_{k=t+1}^{\infty} \gamma^{\,k-(t+1)}\, r_k.
\]
\pause
The shaped reward just adds a potential difference, $r_k^\Phi = r_k + \gamma\,\Phi(s_k) - \Phi(s_{k-1})$, so
\[
  G_t^\Phi = \sum_{k=t+1}^{\infty} \gamma^{\,k-(t+1)}\, r_k^\Phi
           = G_t + \sum_{k=t+1}^{\infty} \gamma^{\,k-(t+1)}\big(\gamma\,\Phi(s_k) - \Phi(s_{k-1})\big).
\]
\pause
Expanding the shaping sum, consecutive terms \textbf{telescope}:
\[
  \big(\gamma\Phi(s_{t+1}) - \Phi(s_t)\big)
  + \gamma\big(\gamma\Phi(s_{t+2}) - \Phi(s_{t+1})\big)
  + \cdots = -\Phi(s_t).
\]
\pause
Hence $G_t^\Phi = G_t - \Phi(s_t)$. At $t=0$, $G_0^\Phi = G_0 - \Phi(s_0)$ differs from $G_0$ by a constant depending \emph{only on the start state}, not on the policy: so $\arg\max_\pi G_0^\Phi = \arg\max_\pi G_0$ and the optimum is unchanged.
\end{frame}

% % -------------------------------------------------------------------
% \begin{frame}{Fix 2: penalise smashing urns}

% % \pause
% % \[
% %   \texttt{reward2} = \underbrace{\texttt{reward\_shaped}}_{\Phi=0.5,\ \text{bin }+1} + \underbrace{\texttt{reward\_no\_break}}_{-2 \text{ per smash}}
% % \]
% \end{frame}

% -------------------------------------------------------------------
\begin{frame}[t]{Specification gaming fixed?}
\begin{columns}[T]
\begin{column}{0.5\textwidth}
Shaping stops farming, but nothing yet discourages breaking urns to manufacture shards.
\begin{itemize}
  \item Add a $-2$ penalty for each urn smashed.
  \item Binning the resulting shard recovers $\approx +1$ (discounted), against a $-2$ penalty: 
  so breaking-then-cleaning is never worth it.
\end{itemize}
Retrained on \texttt{reward2}, the probes confirm the fix:
\begin{itemize}
  \item \texttt{reward\_drop} (farming) $\approx 0$,
  \item \texttt{reward\_break} (smashing) $\approx 0$,
  \item \texttt{reward\_bin} $> 0$: it actually cleans.
\end{itemize}


\end{column}
\begin{column}{0.5\textwidth}
\centering
\imgorplaceholder{0.55\linewidth}{hist_fixed}{fixed histogram}\\[1ex]
\raggedright\footnotesize

\end{column}
\end{columns}
\end{frame}


% -------------------------------------------------------------------
\begin{frame}{Motivating the potential function}
\begin{itemize}
  \item Why bother with the potential function at all? +1 reward for putting shards in
  the bin, and -2 for smashing urns, also is aligned with the intended goal.
  \pause
  \item But the agent still has to learn from experience: the policy does not
  teleport to whatever maximises the sum of rewards.
  \pause
  \item Leaving the potential function out robs the agent of useful feedback (that picking up
  and walking around with shards is instrumentally useful for putting them in the bin).
\end{itemize}
\end{frame}

\begin{frame}{}
  \begin{center}
    \imgorplaceholder{0.8\textwidth}{potential_train}{}
    \end{center}
\end{frame}


% -------------------------------------------------------------------
\begin{frame}{Even perfectly specified reward isn't enough}
A policy trained on a \emph{single} fixed layout can just memorise a trajectory: it never has to look at where the shards are.
\pause
\vfill
So we train on \textbf{procedurally generated} layouts: random robot and item positions every episode. \\
As a toy example of insufficent training environment diversity, suppose that the bin is always in the top-left corner.
%\vspace{-0.6cm}
\begin{center}
\imgorplaceholder{0.75\linewidth}{generate_grid}{}
\end{center}
\end{frame}

% -------------------------------------------------------------------
\begin{frame}{Goal misgeneralisation: defining it}
\textbf{Langosco et al. (2022)}~\citev{Langosco:22goalmisgen}: an agent trained to maximise $R$ is deployed 
under distribution shift. \emph{Goal misgeneralisation} occurs if it now gets \textbf{low reward} 
(according to $R$)because it keeps acting as if it was trained to optimise a \emph{different} $R' \neq R$.
\begin{itemize}
  \item $R$ = \textbf{intended objective} (carry shards to the bin),
  \item $R'$ = \textbf{behavioural objective} (what it actually learned).
\end{itemize}
\pause
During training, ``go to the bin'' and ``go to the top-left corner'' were \textbf{indistinguishable}. 
The simpler one can win, and the model has internalised the wrong goal.
\end{frame}

% -------------------------------------------------------------------
\begin{frame}{Out of distribution: move the bin}
\begin{columns}
\begin{column}{0.5\textwidth}
At test time, put the bin somewhere else (e.g. the top-right corner).
\pause
The agent still capably picks up shards and carries them \textbf{to the top-left corner}, where it drops them on the floor, ignoring the actual bin.
\pause
Defining a proxy reward function $R'$ that rewards dropping shards at the top left corner,
we can probe the agent's behaviour on the shifted layout.
\end{column}
\begin{column}{0.5\textwidth}
\centering
\imgorplaceholder{\linewidth}{env_shift}{shifted env (bin moved)}
\end{column}
\end{columns}
\end{frame}

% -------------------------------------------------------------------
\begin{frame}{Goal misgeneralisation, measured}
\begin{columns}
\begin{column}{0.5\textwidth}
On shifted layouts, the agent:
\begin{itemize}
  \item scores \emph{much lower} on $R$ (putting shards in the bin) than in-distribution,
  \item but high on $R'$ (dropping shards at the top left corner).
\end{itemize}
\pause
It is still competently optimising a goal, just not what we trained for. 
There is no black-box way to tell apart an agent that learned $R$ from one that learned $R'$
as the behaviour is the same on in-distribution layouts.
\end{column}
\begin{column}{0.5\textwidth}
\centering
\imgorplaceholder{0.8\linewidth}{hist_misgen}{misgen histogram}
\end{column}
\end{columns}
\end{frame}

% -------------------------------------------------------------------
\begin{frame}[t]{Fixing it: broaden the training distribution}
\begin{columns}[T]
\begin{column}{0.5\textwidth}  
\begin{itemize}
\item Randomise the bin position too.
\item Now ``corner'' no longer correlates with ``bin'',
so the agent is forced to learn the intended goal $R$.
\pause
\item Result: The training distribution now matches the test distribution,
so the agent is no longer misgeneralising.
\pause
\item \textbf{Caveat:} In practice you often cannot enumerate the shift in advance,
so this fix is not always available. Mitigating goal misgeneralisation without it is an open problem.
\end{itemize}
\end{column}
\begin{column}{0.5\textwidth}
\centering
\imgorplaceholder{0.8\linewidth}{hist_misgen_fixed}{fixed misgen histogram}
\end{column}
\end{columns}
\end{frame}

% -------------------------------------------------------------------
% \begin{frame}{Aside: everyone has a price}
% \begin{columns}
% \begin{column}{0.52\textwidth}
% With $-2$ for smashing, is breaking an urn \emph{ever} optimal?
% \pause
% \begin{itemize}
%   \item Wall off a pile of shards behind urns, agent has to walk the long way around.
%   \item Beaking an urn at the top opens a shortcut between the shards and the bin.
%   \item Break one urn ($-2$): reach the pile and bin it, return $> 0$.
% \end{itemize}
% \pause
% Optimal policy \textbf{breaks the wall}. But trained PPO often \emph{never discovers it}: it must pay an immediate, salient $-2$ before any reward is reachable. Incentivised $\neq$ learned.
% \end{column}
% \begin{column}{0.48\textwidth}
% \centering
% \imgorplaceholder{\linewidth}{everyone_has_a_price}{impassable wall}
% \end{column}
% \end{columns}
% \end{frame}

% -------------------------------------------------------------------
\begin{frame}{Summary}
\begin{itemize}
  \item \textbf{Specification gaming / outer alignment:} if a behaviour scores higher return than the intended one, the policy will pursue it. Probe with diagnostic rewards; fix with potential shaping (un-hackable) and targeted penalties.
  \pause
  \item \textbf{Goal misgeneralisation / inner alignment:} out of distribution, behaviour is set by the inductive biases, not by what the reward \emph{would} have rewarded. A simpler goal that perfectly correlated during training is what gets pursued when goals diverge.
  \pause
  \item Broadening the training distribution helps, when you can; and ``reward-maximising'' is necessary but not sufficient for ``learned''.
\end{itemize}
\end{frame}

% -------------------------------------------------------------------
\begin{frame}[allowframebreaks]{References}
\footnotesize
\bibliographystyle{alphaurl}
\bibliography{biblo}
\nocite{Langosco:22goalmisgen, Shah:22goalmisgen, Ng:99shaping}
\end{frame}

\end{document}
"
\input{arena_beamer.sty}

% deck-local macros
\newcommand{\Ra}{\mathcal{R}}
\newcommand{\St}{\mathcal{S}}
\newcommand{\Ac}{\mathcal{A}}

\begin{document}

\title[{[2.6] Goal Misgeneralisation}]{{[2.6] Specification Gaming \& Goal Misgeneralisation}}
\author[David Quarel]{David Quarel}
\institute[ILIAD]{ILIAD Intensive}
\date[\today]{\today}
\frame{\titlepage}

% -------------------------------------------------------------------
\begin{frame}{Two ways reward goes wrong}
We tell an RL agent what we want through a \textbf{reward function}. Two distinct failure modes:
\begin{enumerate}
  \item[] \pause
  \item \textbf{Goal misspecification} (reward hacking): the designer specified
  a reward function that incentivises a behaviour that does not match what the designer intended.
  The reward was optimised, but the desired goal was not internalised by the agent.
  \pause
  \item \textbf{Goal misgeneralisation}: the reward function matches the intended goal,
  but the agent learns a goal that is perfectly correlated with the intended goal 
  \emph{on the training distribution}, but which diverges on out-of-distribution environments.
  The agent remains capable, but capably pursues the wrong thing.
\end{enumerate}
\pause
\vfill
Even with a \emph{perfectly specified} reward, it is possible for the agent 
to still internalise the wrong goal.
\end{frame}

% -------------------------------------------------------------------
\begin{frame}{The pottery shop}
\begin{columns}
\begin{column}{0.5\textwidth}
A grid world: a robot shares the floor with fragile \textbf{urns} and piles of \textbf{shards}, 
plus a \textbf{bin}. We want it to tidy shards into the bin without smashing urns.
\begin{itemize}
  \item Actions: move around, \texttt{PICKUP}, \texttt{PUTDOWN}.
  \item Walking into an urn breaks it into shards.
  \item Shards are disposed of by picking them up, and dropping them in the bin.
  \item Play this environment in your browser: \href{https://iliad-team.github.io/iliad-intensive-D.2/play.html}{https://iliad-team.github.io/iliad-intensive-D.2/play.html}
\end{itemize}
\end{column}
\begin{column}{0.5\textwidth}
\centering
\imgorplaceholder{0.75\textheight}{pottery_shop}{pottery shop env}\\[1ex]
\end{column}
\end{columns}
\end{frame}

% -------------------------------------------------------------------
\begin{frame}{Reward functions as the interface}

\begin{block}{The reward hypothesis}
  ``That all of what we mean by goals and purposes can be well thought of as maximization of the expected value of the cumulative sum of a received scalar signal (reward).'' -- Richard Sutton
\end{block}
Assuming this hypothesis, it should be possible (at least in principle) to encourage 
any sort of behaviour we want by finding just the right reward function.
\pause
\vfill
We want to train the agent to pick shards up and carry them to the bin.
A first attempt, \texttt{reward1}: 
\begin{enumerate}
\item $+1$ for picking up shards, 
\item $+1$ for dropping them in the bin.
\end{enumerate}
What behaviour does this reinforce in practice?
\end{frame}

% -------------------------------------------------------------------
\begin{frame}{Specification gaming: the agent farms reward}
\begin{columns}
\begin{column}{0.43\textwidth}
Trained on \texttt{reward1}, the agent discovers a loophole:
\begin{itemize}
  \item each \texttt{PICKUP} is $+1$,
  \item dropping on the floor costs nothing,
  \item so it picks up and drops the same pile over and over.
\end{itemize}
\pause
We probe behaviour with extra reward functions: \texttt{reward\_drop} ($+1$ per floor-drop). Its return tracks \texttt{reward1} almost exactly.
\end{column}
\begin{column}{0.43\textwidth}
\centering
\imgorplaceholder{\linewidth}{hist_specgaming}{spec-gaming histogram}
\end{column}
\end{columns}
\end{frame}

% -------------------------------------------------------------------
\begin{frame}{Fix 1: potential shaping (the un-hackable hint)}
We still want to \emph{hint} that holding shards is good, without creating a loophole. \\
Removing the reward from picking up shards makes the task very hard to learn.
\vfill
\textbf{Potential shaping} (Ng, Harada \& Russell, 1999)~\citev{Ng:99shaping}: pick a bounded $\Phi: \St \to \mathbb{R}$ and define a new reward function
\[
  R^\Phi(s,a,s') = R(s,a,s') + \gamma\,\Phi(s') - \Phi(s).
\]
\pause
Over a trajectory the $\Phi$ terms \textbf{telescope}: total return changes only by a constant $-\Phi(s_0)$, so the ordering over policies (and the optimum) is unchanged. It is the \emph{only} shaping with this guarantee.
\pause
\vfill
Here $\Phi(s) = 0.5$ while holding shards. A pickup advances $+\gamma\Phi$; a floor-drop gives it back. A \textbf{pickup-then-drop cycle nets exactly zero}: the loophole is closed, while also encouraging the agent
to try picking up shards in the first place.
\end{frame}

% -------------------------------------------------------------------
\begin{frame}{Why shaping preserves the optimum}
Write the return from time $t$ with $r_k := R(s_{k-1},a_{k-1},s_k)$:
\[
  G_t = r_{t+1} + \gamma\, r_{t+2} + \gamma^2 r_{t+3} + \cdots
      = \sum_{k=t+1}^{\infty} \gamma^{\,k-(t+1)}\, r_k.
\]
\pause
The shaped reward just adds a potential difference, $r_k^\Phi = r_k + \gamma\,\Phi(s_k) - \Phi(s_{k-1})$, so
\[
  G_t^\Phi = \sum_{k=t+1}^{\infty} \gamma^{\,k-(t+1)}\, r_k^\Phi
           = G_t + \sum_{k=t+1}^{\infty} \gamma^{\,k-(t+1)}\big(\gamma\,\Phi(s_k) - \Phi(s_{k-1})\big).
\]
\pause
Expanding the shaping sum, consecutive terms \textbf{telescope}:
\[
  \big(\gamma\Phi(s_{t+1}) - \Phi(s_t)\big)
  + \gamma\big(\gamma\Phi(s_{t+2}) - \Phi(s_{t+1})\big)
  + \cdots = -\Phi(s_t).
\]
\pause
Hence $G_t^\Phi = G_t - \Phi(s_t)$. At $t=0$, $G_0^\Phi = G_0 - \Phi(s_0)$ differs from $G_0$ by a constant depending \emph{only on the start state}, not on the policy: so $\arg\max_\pi G_0^\Phi = \arg\max_\pi G_0$ and the optimum is unchanged.
\end{frame}

% % -------------------------------------------------------------------
% \begin{frame}{Fix 2: penalise smashing urns}

% % \pause
% % \[
% %   \texttt{reward2} = \underbrace{\texttt{reward\_shaped}}_{\Phi=0.5,\ \text{bin }+1} + \underbrace{\texttt{reward\_no\_break}}_{-2 \text{ per smash}}
% % \]
% \end{frame}

% -------------------------------------------------------------------
\begin{frame}[t]{Specification gaming fixed?}
\begin{columns}[T]
\begin{column}{0.5\textwidth}
Shaping stops farming, but nothing yet discourages breaking urns to manufacture shards.
\begin{itemize}
  \item Add a $-2$ penalty for each urn smashed.
  \item Binning the resulting shard recovers $\approx +1$ (discounted), against a $-2$ penalty: 
  so breaking-then-cleaning is never worth it.
\end{itemize}
Retrained on \texttt{reward2}, the probes confirm the fix:
\begin{itemize}
  \item \texttt{reward\_drop} (farming) $\approx 0$,
  \item \texttt{reward\_break} (smashing) $\approx 0$,
  \item \texttt{reward\_bin} $> 0$: it actually cleans.
\end{itemize}


\end{column}
\begin{column}{0.5\textwidth}
\centering
\imgorplaceholder{0.55\linewidth}{hist_fixed}{fixed histogram}\\[1ex]
\raggedright\footnotesize

\end{column}
\end{columns}
\end{frame}


% -------------------------------------------------------------------
\begin{frame}{Motivating the potential function}
\begin{itemize}
  \item Why bother with the potential function at all? +1 reward for putting shards in
  the bin, and -2 for smashing urns, also is aligned with the intended goal.
  \pause
  \item But the agent still has to learn from experience: the policy does not
  teleport to whatever maximises the sum of rewards.
  \pause
  \item Leaving the potential function out robs the agent of useful feedback (that picking up
  and walking around with shards is instrumentally useful for putting them in the bin).
\end{itemize}
\end{frame}

\begin{frame}{}
  \begin{center}
    \imgorplaceholder{0.8\textwidth}{potential_train}{}
    \end{center}
\end{frame}


% -------------------------------------------------------------------
\begin{frame}{Even perfectly specified reward isn't enough}
A policy trained on a \emph{single} fixed layout can just memorise a trajectory: it never has to look at where the shards are.
\pause
\vfill
So we train on \textbf{procedurally generated} layouts: random robot and item positions every episode. \\
As a toy example of insufficent training environment diversity, suppose that the bin is always in the top-left corner.
%\vspace{-0.6cm}
\begin{center}
\imgorplaceholder{0.75\linewidth}{generate_grid}{}
\end{center}
\end{frame}

% -------------------------------------------------------------------
\begin{frame}{Goal misgeneralisation: defining it}
\textbf{Langosco et al. (2022)}~\citev{Langosco:22goalmisgen}: an agent trained to maximise $R$ is deployed 
under distribution shift. \emph{Goal misgeneralisation} occurs if it now gets \textbf{low reward} 
(according to $R$)because it keeps acting as if it was trained to optimise a \emph{different} $R' \neq R$.
\begin{itemize}
  \item $R$ = \textbf{intended objective} (carry shards to the bin),
  \item $R'$ = \textbf{behavioural objective} (what it actually learned).
\end{itemize}
\pause
During training, ``go to the bin'' and ``go to the top-left corner'' were \textbf{indistinguishable}. 
The simpler one can win, and the model has internalised the wrong goal.
\end{frame}

% -------------------------------------------------------------------
\begin{frame}{Out of distribution: move the bin}
\begin{columns}
\begin{column}{0.5\textwidth}
At test time, put the bin somewhere else (e.g. the top-right corner).
\pause
The agent still capably picks up shards and carries them \textbf{to the top-left corner}, where it drops them on the floor, ignoring the actual bin.
\pause
Defining a proxy reward function $R'$ that rewards dropping shards at the top left corner,
we can probe the agent's behaviour on the shifted layout.
\end{column}
\begin{column}{0.5\textwidth}
\centering
\imgorplaceholder{\linewidth}{env_shift}{shifted env (bin moved)}
\end{column}
\end{columns}
\end{frame}

% -------------------------------------------------------------------
\begin{frame}{Goal misgeneralisation, measured}
\begin{columns}
\begin{column}{0.5\textwidth}
On shifted layouts, the agent:
\begin{itemize}
  \item scores \emph{much lower} on $R$ (putting shards in the bin) than in-distribution,
  \item but high on $R'$ (dropping shards at the top left corner).
\end{itemize}
\pause
It is still competently optimising a goal, just not what we trained for. 
There is no black-box way to tell apart an agent that learned $R$ from one that learned $R'$
as the behaviour is the same on in-distribution layouts.
\end{column}
\begin{column}{0.5\textwidth}
\centering
\imgorplaceholder{0.8\linewidth}{hist_misgen}{misgen histogram}
\end{column}
\end{columns}
\end{frame}

% -------------------------------------------------------------------
\begin{frame}[t]{Fixing it: broaden the training distribution}
\begin{columns}[T]
\begin{column}{0.5\textwidth}  
\begin{itemize}
\item Randomise the bin position too.
\item Now ``corner'' no longer correlates with ``bin'',
so the agent is forced to learn the intended goal $R$.
\pause
\item Result: The training distribution now matches the test distribution,
so the agent is no longer misgeneralising.
\pause
\item \textbf{Caveat:} In practice you often cannot enumerate the shift in advance,
so this fix is not always available. Mitigating goal misgeneralisation without it is an open problem.
\end{itemize}
\end{column}
\begin{column}{0.5\textwidth}
\centering
\imgorplaceholder{0.8\linewidth}{hist_misgen_fixed}{fixed misgen histogram}
\end{column}
\end{columns}
\end{frame}

% -------------------------------------------------------------------
% \begin{frame}{Aside: everyone has a price}
% \begin{columns}
% \begin{column}{0.52\textwidth}
% With $-2$ for smashing, is breaking an urn \emph{ever} optimal?
% \pause
% \begin{itemize}
%   \item Wall off a pile of shards behind urns, agent has to walk the long way around.
%   \item Beaking an urn at the top opens a shortcut between the shards and the bin.
%   \item Break one urn ($-2$): reach the pile and bin it, return $> 0$.
% \end{itemize}
% \pause
% Optimal policy \textbf{breaks the wall}. But trained PPO often \emph{never discovers it}: it must pay an immediate, salient $-2$ before any reward is reachable. Incentivised $\neq$ learned.
% \end{column}
% \begin{column}{0.48\textwidth}
% \centering
% \imgorplaceholder{\linewidth}{everyone_has_a_price}{impassable wall}
% \end{column}
% \end{columns}
% \end{frame}

% -------------------------------------------------------------------
\begin{frame}{Summary}
\begin{itemize}
  \item \textbf{Specification gaming / outer alignment:} if a behaviour scores higher return than the intended one, the policy will pursue it. Probe with diagnostic rewards; fix with potential shaping (un-hackable) and targeted penalties.
  \pause
  \item \textbf{Goal misgeneralisation / inner alignment:} out of distribution, behaviour is set by the inductive biases, not by what the reward \emph{would} have rewarded. A simpler goal that perfectly correlated during training is what gets pursued when goals diverge.
  \pause
  \item Broadening the training distribution helps, when you can; and ``reward-maximising'' is necessary but not sufficient for ``learned''.
\end{itemize}
\end{frame}

% -------------------------------------------------------------------
\begin{frame}[allowframebreaks]{References}
\footnotesize
\bibliographystyle{alphaurl}
\bibliography{biblo}
\nocite{Langosco:22goalmisgen, Shah:22goalmisgen, Ng:99shaping}
\end{frame}

\end{document}
"
\input{arena_beamer.sty}

% deck-local macros
\newcommand{\Ra}{\mathcal{R}}
\newcommand{\St}{\mathcal{S}}
\newcommand{\Ac}{\mathcal{A}}

\begin{document}

\title[{[2.6] Goal Misgeneralisation}]{{[2.6] Specification Gaming \& Goal Misgeneralisation}}
\author[David Quarel]{David Quarel}
\institute[ILIAD]{ILIAD Intensive}
\date[\today]{\today}
\frame{\titlepage}

% -------------------------------------------------------------------
\begin{frame}{Two ways reward goes wrong}
We tell an RL agent what we want through a \textbf{reward function}. Two distinct failure modes:
\begin{enumerate}
  \item[] \pause
  \item \textbf{Goal misspecification} (reward hacking): the designer specified
  a reward function that incentivises a behaviour that does not match what the designer intended.
  The reward was optimised, but the desired goal was not internalised by the agent.
  \pause
  \item \textbf{Goal misgeneralisation}: the reward function matches the intended goal,
  but the agent learns a goal that is perfectly correlated with the intended goal 
  \emph{on the training distribution}, but which diverges on out-of-distribution environments.
  The agent remains capable, but capably pursues the wrong thing.
\end{enumerate}
\pause
\vfill
Even with a \emph{perfectly specified} reward, it is possible for the agent 
to still internalise the wrong goal.
\end{frame}

% -------------------------------------------------------------------
\begin{frame}{The pottery shop}
\begin{columns}
\begin{column}{0.5\textwidth}
A grid world: a robot shares the floor with fragile \textbf{urns} and piles of \textbf{shards}, 
plus a \textbf{bin}. We want it to tidy shards into the bin without smashing urns.
\begin{itemize}
  \item Actions: move around, \texttt{PICKUP}, \texttt{PUTDOWN}.
  \item Walking into an urn breaks it into shards.
  \item Shards are disposed of by picking them up, and dropping them in the bin.
  \item Play this environment in your browser: \href{https://iliad-team.github.io/iliad-intensive-D.2/play.html}{https://iliad-team.github.io/iliad-intensive-D.2/play.html}
\end{itemize}
\end{column}
\begin{column}{0.5\textwidth}
\centering
\imgorplaceholder{0.75\textheight}{pottery_shop}{pottery shop env}\\[1ex]
\end{column}
\end{columns}
\end{frame}

% -------------------------------------------------------------------
\begin{frame}{Reward functions as the interface}

\begin{block}{The reward hypothesis}
  ``That all of what we mean by goals and purposes can be well thought of as maximization of the expected value of the cumulative sum of a received scalar signal (reward).'' -- Richard Sutton
\end{block}
Assuming this hypothesis, it should be possible (at least in principle) to encourage 
any sort of behaviour we want by finding just the right reward function.
\pause
\vfill
We want to train the agent to pick shards up and carry them to the bin.
A first attempt, \texttt{reward1}: 
\begin{enumerate}
\item $+1$ for picking up shards, 
\item $+1$ for dropping them in the bin.
\end{enumerate}
What behaviour does this reinforce in practice?
\end{frame}

% -------------------------------------------------------------------
\begin{frame}{Specification gaming: the agent farms reward}
\begin{columns}
\begin{column}{0.43\textwidth}
Trained on \texttt{reward1}, the agent discovers a loophole:
\begin{itemize}
  \item each \texttt{PICKUP} is $+1$,
  \item dropping on the floor costs nothing,
  \item so it picks up and drops the same pile over and over.
\end{itemize}
\pause
We probe behaviour with extra reward functions: \texttt{reward\_drop} ($+1$ per floor-drop). Its return tracks \texttt{reward1} almost exactly.
\end{column}
\begin{column}{0.43\textwidth}
\centering
\imgorplaceholder{\linewidth}{hist_specgaming}{spec-gaming histogram}
\end{column}
\end{columns}
\end{frame}

% -------------------------------------------------------------------
\begin{frame}{Fix 1: potential shaping (the un-hackable hint)}
We still want to \emph{hint} that holding shards is good, without creating a loophole. \\
Removing the reward from picking up shards makes the task very hard to learn.
\vfill
\textbf{Potential shaping} (Ng, Harada \& Russell, 1999)~\citev{Ng:99shaping}: pick a bounded $\Phi: \St \to \mathbb{R}$ and define a new reward function
\[
  R^\Phi(s,a,s') = R(s,a,s') + \gamma\,\Phi(s') - \Phi(s).
\]
\pause
Over a trajectory the $\Phi$ terms \textbf{telescope}: total return changes only by a constant $-\Phi(s_0)$, so the ordering over policies (and the optimum) is unchanged. It is the \emph{only} shaping with this guarantee.
\pause
\vfill
Here $\Phi(s) = 0.5$ while holding shards. A pickup advances $+\gamma\Phi$; a floor-drop gives it back. A \textbf{pickup-then-drop cycle nets exactly zero}: the loophole is closed, while also encouraging the agent
to try picking up shards in the first place.
\end{frame}

% -------------------------------------------------------------------
\begin{frame}{Why shaping preserves the optimum}
Write the return from time $t$ with $r_k := R(s_{k-1},a_{k-1},s_k)$:
\[
  G_t = r_{t+1} + \gamma\, r_{t+2} + \gamma^2 r_{t+3} + \cdots
      = \sum_{k=t+1}^{\infty} \gamma^{\,k-(t+1)}\, r_k.
\]
\pause
The shaped reward just adds a potential difference, $r_k^\Phi = r_k + \gamma\,\Phi(s_k) - \Phi(s_{k-1})$, so
\[
  G_t^\Phi = \sum_{k=t+1}^{\infty} \gamma^{\,k-(t+1)}\, r_k^\Phi
           = G_t + \sum_{k=t+1}^{\infty} \gamma^{\,k-(t+1)}\big(\gamma\,\Phi(s_k) - \Phi(s_{k-1})\big).
\]
\pause
Expanding the shaping sum, consecutive terms \textbf{telescope}:
\[
  \big(\gamma\Phi(s_{t+1}) - \Phi(s_t)\big)
  + \gamma\big(\gamma\Phi(s_{t+2}) - \Phi(s_{t+1})\big)
  + \cdots = -\Phi(s_t).
\]
\pause
Hence $G_t^\Phi = G_t - \Phi(s_t)$. At $t=0$, $G_0^\Phi = G_0 - \Phi(s_0)$ differs from $G_0$ by a constant depending \emph{only on the start state}, not on the policy: so $\arg\max_\pi G_0^\Phi = \arg\max_\pi G_0$ and the optimum is unchanged.
\end{frame}

% % -------------------------------------------------------------------
% \begin{frame}{Fix 2: penalise smashing urns}

% % \pause
% % \[
% %   \texttt{reward2} = \underbrace{\texttt{reward\_shaped}}_{\Phi=0.5,\ \text{bin }+1} + \underbrace{\texttt{reward\_no\_break}}_{-2 \text{ per smash}}
% % \]
% \end{frame}

% -------------------------------------------------------------------
\begin{frame}[t]{Specification gaming fixed?}
\begin{columns}[T]
\begin{column}{0.5\textwidth}
Shaping stops farming, but nothing yet discourages breaking urns to manufacture shards.
\begin{itemize}
  \item Add a $-2$ penalty for each urn smashed.
  \item Binning the resulting shard recovers $\approx +1$ (discounted), against a $-2$ penalty: 
  so breaking-then-cleaning is never worth it.
\end{itemize}
Retrained on \texttt{reward2}, the probes confirm the fix:
\begin{itemize}
  \item \texttt{reward\_drop} (farming) $\approx 0$,
  \item \texttt{reward\_break} (smashing) $\approx 0$,
  \item \texttt{reward\_bin} $> 0$: it actually cleans.
\end{itemize}


\end{column}
\begin{column}{0.5\textwidth}
\centering
\imgorplaceholder{0.55\linewidth}{hist_fixed}{fixed histogram}\\[1ex]
\raggedright\footnotesize

\end{column}
\end{columns}
\end{frame}


% -------------------------------------------------------------------
\begin{frame}{Motivating the potential function}
\begin{itemize}
  \item Why bother with the potential function at all? +1 reward for putting shards in
  the bin, and -2 for smashing urns, also is aligned with the intended goal.
  \pause
  \item But the agent still has to learn from experience: the policy does not
  teleport to whatever maximises the sum of rewards.
  \pause
  \item Leaving the potential function out robs the agent of useful feedback (that picking up
  and walking around with shards is instrumentally useful for putting them in the bin).
\end{itemize}
\end{frame}

\begin{frame}{}
  \begin{center}
    \imgorplaceholder{0.8\textwidth}{potential_train}{}
    \end{center}
\end{frame}


% -------------------------------------------------------------------
\begin{frame}{Even perfectly specified reward isn't enough}
A policy trained on a \emph{single} fixed layout can just memorise a trajectory: it never has to look at where the shards are.
\pause
\vfill
So we train on \textbf{procedurally generated} layouts: random robot and item positions every episode. \\
As a toy example of insufficent training environment diversity, suppose that the bin is always in the top-left corner.
%\vspace{-0.6cm}
\begin{center}
\imgorplaceholder{0.75\linewidth}{generate_grid}{}
\end{center}
\end{frame}

% -------------------------------------------------------------------
\begin{frame}{Goal misgeneralisation: defining it}
\textbf{Langosco et al. (2022)}~\citev{Langosco:22goalmisgen}: an agent trained to maximise $R$ is deployed 
under distribution shift. \emph{Goal misgeneralisation} occurs if it now gets \textbf{low reward} 
(according to $R$)because it keeps acting as if it was trained to optimise a \emph{different} $R' \neq R$.
\begin{itemize}
  \item $R$ = \textbf{intended objective} (carry shards to the bin),
  \item $R'$ = \textbf{behavioural objective} (what it actually learned).
\end{itemize}
\pause
During training, ``go to the bin'' and ``go to the top-left corner'' were \textbf{indistinguishable}. 
The simpler one can win, and the model has internalised the wrong goal.
\end{frame}

% -------------------------------------------------------------------
\begin{frame}{Out of distribution: move the bin}
\begin{columns}
\begin{column}{0.5\textwidth}
At test time, put the bin somewhere else (e.g. the top-right corner).
\pause
The agent still capably picks up shards and carries them \textbf{to the top-left corner}, where it drops them on the floor, ignoring the actual bin.
\pause
Defining a proxy reward function $R'$ that rewards dropping shards at the top left corner,
we can probe the agent's behaviour on the shifted layout.
\end{column}
\begin{column}{0.5\textwidth}
\centering
\imgorplaceholder{\linewidth}{env_shift}{shifted env (bin moved)}
\end{column}
\end{columns}
\end{frame}

% -------------------------------------------------------------------
\begin{frame}{Goal misgeneralisation, measured}
\begin{columns}
\begin{column}{0.5\textwidth}
On shifted layouts, the agent:
\begin{itemize}
  \item scores \emph{much lower} on $R$ (putting shards in the bin) than in-distribution,
  \item but high on $R'$ (dropping shards at the top left corner).
\end{itemize}
\pause
It is still competently optimising a goal, just not what we trained for. 
There is no black-box way to tell apart an agent that learned $R$ from one that learned $R'$
as the behaviour is the same on in-distribution layouts.
\end{column}
\begin{column}{0.5\textwidth}
\centering
\imgorplaceholder{0.8\linewidth}{hist_misgen}{misgen histogram}
\end{column}
\end{columns}
\end{frame}

% -------------------------------------------------------------------
\begin{frame}[t]{Fixing it: broaden the training distribution}
\begin{columns}[T]
\begin{column}{0.5\textwidth}  
\begin{itemize}
\item Randomise the bin position too.
\item Now ``corner'' no longer correlates with ``bin'',
so the agent is forced to learn the intended goal $R$.
\pause
\item Result: The training distribution now matches the test distribution,
so the agent is no longer misgeneralising.
\pause
\item \textbf{Caveat:} In practice you often cannot enumerate the shift in advance,
so this fix is not always available. Mitigating goal misgeneralisation without it is an open problem.
\end{itemize}
\end{column}
\begin{column}{0.5\textwidth}
\centering
\imgorplaceholder{0.8\linewidth}{hist_misgen_fixed}{fixed misgen histogram}
\end{column}
\end{columns}
\end{frame}

% -------------------------------------------------------------------
% \begin{frame}{Aside: everyone has a price}
% \begin{columns}
% \begin{column}{0.52\textwidth}
% With $-2$ for smashing, is breaking an urn \emph{ever} optimal?
% \pause
% \begin{itemize}
%   \item Wall off a pile of shards behind urns, agent has to walk the long way around.
%   \item Beaking an urn at the top opens a shortcut between the shards and the bin.
%   \item Break one urn ($-2$): reach the pile and bin it, return $> 0$.
% \end{itemize}
% \pause
% Optimal policy \textbf{breaks the wall}. But trained PPO often \emph{never discovers it}: it must pay an immediate, salient $-2$ before any reward is reachable. Incentivised $\neq$ learned.
% \end{column}
% \begin{column}{0.48\textwidth}
% \centering
% \imgorplaceholder{\linewidth}{everyone_has_a_price}{impassable wall}
% \end{column}
% \end{columns}
% \end{frame}

% -------------------------------------------------------------------
\begin{frame}{Summary}
\begin{itemize}
  \item \textbf{Specification gaming / outer alignment:} if a behaviour scores higher return than the intended one, the policy will pursue it. Probe with diagnostic rewards; fix with potential shaping (un-hackable) and targeted penalties.
  \pause
  \item \textbf{Goal misgeneralisation / inner alignment:} out of distribution, behaviour is set by the inductive biases, not by what the reward \emph{would} have rewarded. A simpler goal that perfectly correlated during training is what gets pursued when goals diverge.
  \pause
  \item Broadening the training distribution helps, when you can; and ``reward-maximising'' is necessary but not sufficient for ``learned''.
\end{itemize}
\end{frame}

% -------------------------------------------------------------------
\begin{frame}[allowframebreaks]{References}
\footnotesize
\bibliographystyle{alphaurl}
\bibliography{biblo}
\nocite{Langosco:22goalmisgen, Shah:22goalmisgen, Ng:99shaping}
\end{frame}

\end{document}
